\chapter{Physical questions and comparison conventions}
\label{narr:scope}

The central question is how much of a quantum spin system can be described by a
self-consistent Gaussian state, and which physical predictions survive comparison
with exact finite systems. The Kitaev honeycomb model provides a useful testing
ground: its zero-field integrable structure exposes distinctions between spin
states, auxiliary fermions, and flux sectors that an approximation can obscure.
A magnetic field makes these distinctions more consequential, because it couples
sectors that were conserved at zero field.

This book separates the argument a reader needs from the historical route by
which it was developed. The main chapters establish conventions, finite-state
identities, and the requirements for interpreting a calculation. Technical
appendices give their proofs and preserve older derivations and experiments.
An archived assertion is a record of a research decision, not a renewed theorem.
The current acceptance register takes precedence over old words such as
``certified'' in those records.

% EDITORIAL QUESTION: Which Hamiltonian and parameter domain are discussed here?
\section{The finite physical model}
\label{narr:model}
For a finite honeycomb graph with $N$ sites, take one Pauli algebra at each site
and count each undirected bond once. The reference Hamiltonian is
\begin{equation}
 H(K,\bm h)=-K\sum_{\langle ij\rangle_\alpha}
 \sigma_i^\alpha\sigma_j^\alpha-\sum_i\bm h\cdot\bm\sigma_i,
 \qquad K>0,\qquad \bm h=\frac{h}{\sqrt3}(1,1,1).
 \label{eq:narr:model}
\end{equation}
Here $K>0$ means ferromagnetic exchange, and $h$ is the Euclidean field magnitude,
not one Cartesian component. We set $k_B=\hbar=1$. A periodic sample needs its
primitive vectors, bond list, sizes, and boundary identifications specified in
addition to $N$. An isolated hexagon with missing exterior bonds is a different
Hamiltonian from a honeycomb torus with the same number of sites.

Equation~\eqref{eq:narr:model} fixes the convention for this narrative. It does
not silently identify every inherited parameter named \texttt{K} with $K$.
If a program implements $+K_{\rm code}\sum\sigma_i^\alpha\sigma_j^\alpha$ using
the same spin operators and bonds, equality of Hamiltonians requires
$K_{\rm code}=-K$. A fermionic implementation needs its operator map checked as
well; a parameter name cannot establish that equality.

The present research concerns the pure ferromagnetic model and a [111] field.
Antiferromagnetic exchange, additional couplings, other field directions, and
other lattices introduce new physical questions. The finite-state foundations
below can be reused, but a measured phase boundary or excitation coefficient
cannot be transferred by changing a label.

% EDITORIAL QUESTION: What does changing from Pauli operators to spin-half operators actually change?
\section{Units before numerical comparison}
% CLAIM: W001-FND-06 (note: doc/audits/w001_foundation_review.md).
If $S_i^\alpha=\sigma_i^\alpha/2$, the same operator is written
\begin{equation}
 H=-J\sum_{\langle ij\rangle_\alpha} S_i^\alpha S_j^\alpha
       -\sum_i\bm h_S\cdot\bm S_i,
 \qquad J=4K,\qquad \bm h_S=2\bm h.
 \label{eq:narr:units}
\end{equation}
The energy spectrum is unchanged when the parameters and operators are changed
together. Numerical coefficients quoted in different energy units can change;
that is a separate substitution. In particular, Eq.~\eqref{eq:narr:units} by
itself says nothing about whether a plotted quantity is a positive fermion
energy, a separation between two matrix bands, or a physical spin excitation.
Those observables are distinguished in Chapter~\ref{narr:comparisons}.

A meaningful comparison therefore carries a short dictionary: Hamiltonian,
geometry, state space, field components, energy units, observable, and limiting
procedure. This is the scientific content of a convention map. Agreement of one
number after a rescaling is weaker evidence than agreement of the mapped
operators and several independent observables.

% EDITORIAL QUESTION: How should the reader interpret the evidence in this edition?
\section{What is established and what remains research}
The new logical narrative uses exact finite-dimensional identities with explicit
proofs in Appendix~\ref{app:finite_foundations}. Conditional perturbative results
retain their invertibility and norm hypotheses. Numerical values inherited from
older waves appear in historical appendices with their current review status.
They are not used here to certify a phase diagram or an exact prefactor.

The original perturbation framework is an essential prerequisite: it distinguishes
the physical flux-free projector, a ground-eigenspace projector, an external
energy window, and a one-particle occupied projector. These distinctions were
ignored in some later interpretations; the records have been preserved rather
than rewritten to conceal that problem. The historical analytical/MF-versus-ED
study called Phase 4 is disapproved by the principal investigator. Its appendix
is retained so that its assumptions can be inspected, not used as an accepted
explanation of a discrepancy.
